\section{Electron trajectory, return spectrum, and mass conservation}
\label{vi:sec:trayectoria-espectral-masa}

This section develops, on the mass operator already constructed, the
dynamical structure that preserves the electron trajectory. It introduces
neither a second definition of the electron nor an additional calibration.
The causal chain is the one established throughout this article:
\[
 \mathrm{APP}\longrightarrow\mathrm{TRIT}\longrightarrow\mathrm{TPK}
 \longrightarrow X_{\rm enr}
 \longrightarrow\mathcal C_{\rm cont}^{\rm disc}
 \longrightarrow (\Gamma_e,\tau_e,\sigma_e)
 \longrightarrow M^\beta .
\]
The coordinates $\pi_{\rm HMT},\varphi_{\rm HMT},e_{\rm HMT}$, and
$\alpha_{\rm HMT}$ are correlated HMT outputs generated from the enriched
state; they are not introduced as conventional input data. The Fourier analysis
below is applied after the word has been produced and does not select the route.

\subsection{Distinct central objects and fibre memory}

The fixed cell $(5,5)$, the boundary datum $111111$, the transverse emission
$555555$, and the orientation $\tau_e=(+++---)^2$ do not name one and the
same object. The cellular involution
\[
 \rho_c(r,c)=(10-r,10-c)
\]
has the unique fixed point $(5,5)$. At that point,
\[
 5+5=1+9,\qquad 5\cdot5=7+2\cdot9,
\]
whereas at $(8,2)$ and $(2,8)$ the sum is preserved but the product quotient
changes. The fixed cell determines the central exception; the boundary
register and the transverse emission preserve other state coordinates.

In particular, the product-channel fibre of $555555$ contains $24$ oriented
positions. After one TPK advance it splits into
\[
 555177,\qquad 555717,\qquad 555771,
\]
with eight positions in each output. The next value is not a function of the
isolated scalar signature: the memory class and orientation are required to
reconstruct the evolution. This is why the mass operator acts on route fibres
rather than on a list of numbers.

\subsection{Exact cellular support of the central trajectory}

The trajectory in~\eqref{vi:dep:trayectoria-electron} starts at
$(r,c,h,v)=(5,5,1,1)$. Each TRIT block $(+1,-1,0)$ advances each coordinate
once and retains one threshold step. At the endpoints of every block one has
$r=c$; at the intermediate step,
\[
 c-r\equiv h\pmod 9.
\]
Nine blocks visit every position on the corresponding diagonal. The
simultaneous inversions of $(h,v)$ at steps $27,54,81$ add the conjugate
diagonal without resetting the state. Consequently, the cellular projection
of the $108$ steps has the exact support
\begin{equation}
 \operatorname{Supp}_{\rm cell}(\Gamma_e)
 =\left\{(r,c)\in\{1,\ldots,9\}^2:
 c-r\equiv0,1,-1\pmod9\right\},
 \qquad
 \left|\operatorname{Supp}_{\rm cell}(\Gamma_e)\right|=27.
 \label{vi:eq:soporte-electron-27}
\end{equation}
The equality follows from the update rule, not from a retrospective count.
It describes a discrete trajectory morphology on the positional torus. Its
realisation as a M\"obius band additionally requires the orientation transport
and gluing of Section~\ref{vi:sec:conjugacion}; the finite support is not
identified with a three-dimensional material shape.

\subsection{Two histories, one sector chart, and exact reconstruction}

The central biography contains the histories
\begin{equation}
 e_{++}=(2,2,5,-4,-3,-2)^2,
 \qquad
 e_{--}=(4,3,2,-2,-2,-5)^2.
 \label{vi:eq:historias-electronicas}
\end{equation}
Both are expressed in the fixed sector chart
$\Sigma_{12}=(+++---)^2$. The absolute orientation of the recurrence is
transported separately; the sign of a history is not identified with the
sector label. The two histories have zero total and the same absolute-value
histogram, but they do not have the same positional memory.

For a history $e=(e_0,\ldots,e_{11})$, define
\begin{equation}
 p_i=e_i+e_{i+6},\qquad a_i=e_i-e_{i+6},
 \qquad s_C=\sum_{i=0}^{5}p_i,
 \qquad M_i=p_i-p_5\quad(0\leq i<5).
 \label{vi:eq:codificador-memoria-electron}
\end{equation}
The decoder is
\begin{equation}
 p_5=\frac{s_C-\sum_{i=0}^{4}M_i}{6},\quad
 p_i=p_5+M_i,\quad
 e_i=\frac{p_i+a_i}{2},\quad e_{i+6}=\frac{p_i-a_i}{2}.
 \label{vi:eq:decodificador-memoria-electron}
\end{equation}
The integrality conditions are divisibility by six in the first expression
and common parity of $p_i$ and $a_i$. For
\eqref{vi:eq:historias-electronicas}, one obtains
\[
 M^{++}=(8,8,14,-4,-2),\qquad
 M^{--}=(18,16,14,6,6),\qquad a=0.
\]
Substitution in~\eqref{vi:eq:decodificador-memoria-electron} restores all twelve
components. Thus the histories are demonstrably distinct even when a partial
reading coincides. It does not follow that their masses must differ: a mass
reader may be invariant under a transport that preserves fewer coordinates
than the enriched state.

\subsection{Spectral identity of the return reader}

Let $\tau=(\tau_0,\ldots,\tau_{11})\in\{-1,+1\}^{12}$ with cyclic indices.
Define
\begin{equation}
 C_s(\tau)=\sum_{j=0}^{11}\tau_j\tau_{j+s},
 \qquad
 \Omega(\tau)=-2C_1-4C_2-6C_3,
 \qquad
 P_6(\tau)=\frac{C_6}{2}.
 \label{vi:eq:lectores-correlacion}
\end{equation}
The first identity is the expanded form of $\Omega=4A_3-3A_4$ already
proved in~\eqref{vi:dep:KOmega}. To make explicit which modes transport the
correction, introduce, after $\tau$ has been produced,
\begin{equation}
 T_k(\tau)=\sum_{j=0}^{11}\tau_j
 e^{-2\pi i k j/12},\qquad 0\leq k<12.
 \label{vi:eq:dft-retorno}
\end{equation}

\begin{proposition}[Exact spectral decomposition]
For every real word of length twelve,
\begin{align}
 C_s&=\frac1{12}\sum_{k=0}^{11}|T_k|^2
          \cos\!\left(\frac{2\pi ks}{12}\right),\label{vi:eq:wiener-khinchin-finito}\\
 \Omega&=\frac1{12}\sum_{k=0}^{11}|T_k|^2w_k,
 \quad
 w_k=-2\cos\theta_k-4\cos2\theta_k-6\cos3\theta_k,
 \quad\theta_k=\frac{2\pi k}{12},\label{vi:eq:omega-espectral}\\
 P_6&=\frac1{24}\sum_{k=0}^{11}|T_k|^2(-1)^k.
 \label{vi:eq:p6-espectral}
\end{align}
\end{proposition}
\begin{proof}
Expanding $|T_k|^2$ yields
$\sum_{j,\ell}\tau_j\tau_\ell e^{-2\pi ik(j-\ell)/12}$.
The sum over $k$ equals twelve when $j-\ell$ is the prescribed displacement
and vanishes otherwise. This is finite character orthogonality and proves
\eqref{vi:eq:wiener-khinchin-finito}. Inserting successively $s=1,2,3$
in~\eqref{vi:eq:lectores-correlacion} yields
\eqref{vi:eq:omega-espectral}; $s=6$ yields
\eqref{vi:eq:p6-espectral}.
\end{proof}

For $\tau_e=(+++---)^2$, with $z_k=e^{-2\pi ik/12}$,
\begin{equation}
 T_k=(1+(-1)^k)(1-z_k^3)(1+z_k+z_k^2).
 \label{vi:eq:factorizacion-tk-electron}
\end{equation}
Only $k=2,6,10$ survive, with
\[
 T_2=4-4i\sqrt3,\qquad T_6=4,\qquad
 T_{10}=4+4i\sqrt3,
\]
and hence
\begin{equation}
 |T_2|^2=64,\qquad |T_6|^2=16,\qquad |T_{10}|^2=64.
 \label{vi:eq:potencias-espectrales-electron}
\end{equation}
The corresponding weights $w_2,w_6,w_{10}$ are $7,4,7$. Therefore,
\begin{equation}
 \Omega(\tau_e)=\frac{64\cdot7+16\cdot4+64\cdot7}{12}=80,
 \qquad P_6(\tau_e)=6.
 \label{vi:eq:omega-p6-electron}
\end{equation}
This recovers the second reader of~\eqref{vi:dep:KOmega} from the trajectory
modes. The alternating word $(+-)^6$, which retains the same sign histogram,
has $|T_6|^2=144$ and gives $\Omega=48$; it is a negative control showing
that order, not merely counting, enters. Cyclic translations multiply $T_k$
by a phase and preserve $|T_k|^2$, so this reader alone does not recover the
time origin.

The return contribution to the electron exponent is thus
\begin{equation}
 \kappa_e=
 \frac1{12}\sum_{k=0}^{11}|T_k|^2
 \left(w_k+\frac{\Delta_4}{2}(-1)^k\right)-54\alpha_{\rm HMT}
 =80-54\alpha_{\rm HMT}+6\Delta_4.
 \label{vi:eq:kappa-espectral-electron}
\end{equation}
This identity introduces no spatial Hamiltonian: it proves that the
already-defined mass correction is an exact spectral functional of the return
word.

\subsection{Extension of the electron equation to multisections}

The register of a route $\gamma$ supplies
$\sigma_\gamma=(n_A,s_C,k_\Delta,b_{120},b_\zeta)$ and tower coordinates
$\eta_\gamma=(\eta_h)_h$. In a common chart, the refined character is
\begin{equation}
 \begin{aligned}
 \log\chi_{\rm ref}(\sigma_\gamma,\eta_\gamma)
 &=n_Ax+\frac{s_C}{6}y+k_\Delta\Delta_4+b_\zeta\zeta_{270}
   +b_{120}s_{120}+\sum_{h\in\langle90,120\rangle}\eta_hs_h\\
 &=n_Ax+\frac{s_C}{6}y+k_\Delta\Delta_4+b_\zeta\zeta_{270}
   +\sum_{h\in\langle90,120\rangle}N_hs_h,\\
 N_h&:=\eta_h+b_{120}\mathbf1_{\{h=120\}}.
 \end{aligned}
 \label{vi:eq:caracter-refinado-explicito}
\end{equation}
where $x=A_{\rm rad}$, $y=C^*_{\rm rad}$,
$\zeta_{270}=135\Delta_4^2$, and $s_h=q_-^h-q_+^h$. The two lines are the
same evaluation: $b_{120}$ occurs once, either as an explicit term or absorbed
into $N_{120}$. Absolute convergence of the last series is required for an
infinite tower. This formula keeps separate the torsion $\zeta_{270}$, the two
sheets, and the two compositions of height $360=4\cdot90=3\cdot120$.

In the route basis of a fibre, let $B^r e_\gamma=\kappa_\gamma^r e_\gamma$,
with $r=D$ or $\Omega$. The operator law of
Section~\ref{vi:sec:operador-masa} gives, in its diagonal chart,
\begin{equation}
 \boxed{
 \frac{m_\gamma^\beta}{m_e^\beta}=
 \chi_{\rm ref}(\sigma_\gamma-\sigma_e,
                 \eta_\gamma-\eta_e)
 R_{\rm act}^{\kappa_\gamma-\kappa_e}.}
 \label{vi:eq:ley-masas-trayectoria}
\end{equation}
The electron equation is the reference section of this identity, not an
experimental number used to select routes. Physical labels are assigned only
after fibres, representations, and channels have been constructed; numerical
proximity is not a substitute for that prospective matching.

\subsection{Operator criterion for mass-preserving transport}

Define the dimensionless coordinate
\begin{equation}
 \Lambda_\gamma:=\log\frac{m_\gamma^\beta}{m_e^\beta}
 \label{vi:eq:lambda-masa-ruta}
\end{equation}
in a common chart with positive masses. This symbol is neither the reciprocal
length $\Lambda=L_\alpha X^{-1}$ nor the coordinate
$\bar\lambda_\Gamma$ used elsewhere.

Let $\mathscr D$ be a finite route domain stable under a bijective TPK
transport $F:\mathscr D\to\mathscr D$, and let
$V_Fe_\gamma=e_{F\gamma}$ be its permutation operator. Then
\begin{equation}
 [M^\beta,V_F]e_\gamma
 =m_e^\beta\left(e^{\Lambda_{F\gamma}}-e^{\Lambda_\gamma}\right)
 e_{F\gamma}.
 \label{vi:eq:conmutador-masa-transporte}
\end{equation}
Consequently,
\begin{equation}
 [M^\beta,V_F]=0
 \quad\Longleftrightarrow\quad
 \Lambda_{F\gamma}=\Lambda_\gamma
 \quad\text{for every }\gamma\in\mathscr D.
 \label{vi:eq:criterio-conservacion-masa}
\end{equation}
Using~\eqref{vi:eq:caracter-refinado-explicito}, this condition is
\begin{equation}
 \delta n_Ax+\frac{\delta s_C}{6}y
 +\delta k_\Delta\Delta_4+\delta b_\zeta\zeta_{270}
 +\sum_h\delta N_hs_h
 +\delta\kappa\log R_{\rm act}=0.
 \label{vi:eq:balance-conservacion-masa}
\end{equation}
Here
\(
 \delta N_h=\delta\eta_h+
 \delta b_{120}\mathbf1_{\{h=120\}},
\)
so the difference retains the original tower coordinates and adds the
$b_{120}$ incidence term exactly once.
It is checked after $F$ has been produced by the TPK; $F$ is not selected
retrospectively by the mass value. The criterion permits events and memory to
change provided the total coordinate is preserved. It does not state that
every TPK transport preserves mass, nor that it exhausts bound-state dynamics.
Its hypotheses are a stable domain, bijective transport, a fixed chart,
positive masses, and convergence of the towers.

\subsection{Joint result}

Equations~\eqref{vi:eq:soporte-electron-27},
\eqref{vi:eq:decodificador-memoria-electron},
\eqref{vi:eq:kappa-espectral-electron},
\eqref{vi:eq:ley-masas-trayectoria}, and
\eqref{vi:eq:criterio-conservacion-masa} close one chain. The central
trajectory visits exactly twenty-seven cells; its histories preserve
reversible memory; the return has three nonzero modes producing
$(\Omega,P_6)=(80,6)$; the mass law extends the electron section to the
multisections; and the commutator identifies exactly which transports
preserve mass. Cellular morphology, return spectrum, and magnitude are
coordinated readers of one enriched state, not interchangeable observables.
